\documentclass{situationist-preprint}

\title{Profound Boredom as a Limit Case}
\msid{PHE-04-b}
\msfamily{PHE-04}
\msdiscipline{Phenomenology}
\mstype{Consequence}
\mskeywords{intentionality; profound boredom; Heidegger; time; moods}

\begin{document}
\maketitle

\begin{abstract}
PHE-04-a argues that consciousness is always consciousness of something, and that profound boredom, the apparent exception, is intentional towards time itself. We derive the consequence that profound boredom is a limit case of intentionality in a precise sense. Heidegger's three forms of boredom form a series in which the object widens while the aspect stays constant: a task, a situation, the time of one's existence. We show that the series has a last member, that this member is still intentional, and that no conscious state lies beyond it in the same series. The limit is reached without being exceeded, so it confirms rather than threatens the intentionality thesis.
\end{abstract}

\section{The series of boredom}

In the lecture course of 1929--30, Heidegger distinguishes three forms of boredom \citep{Heidegger1983}. In the first, one is bored \emph{by} something: waiting at a provincial station for a train, one finds the time dragging and the place offering nothing. In the second, one is bored \emph{with} something: after an evening among friends in which nothing was wrong, one finds that the whole evening was boring, though no item in it was. In the third, profound boredom, \enquote{it is boring for one}: nothing in particular bores, no situation bores, and beings as a whole withdraw into indifference.

PHE-04-a uses Definition~2.1 (a state is intentional if it presents something as being some way) and argues, in Lemma~4.3, that what profound boredom presents as dragging is time, the time of the subject's existence as such; Theorem~4.4 concludes that consciousness is intentional without exception. We make explicit the structure that the three forms share.

\begin{definition}\label{d:series}
The \emph{boredom series} is the sequence of the three forms, ordered so that the object of each later form includes the object of each earlier one as a part, while the aspect, dragging and emptiness, stays the same.
\end{definition}

The object of the first form is a determinate stretch of time occupied by a task or a place; of the second, a situation as a whole, an evening, in which the time one has given is presented as passed emptily; of the third, by Lemma~4.3 of PHE-04-a, the time of one's existence. Each object includes the earlier ones: the evening contains the waiting, and one's existence contains every evening.

\section{The limit}

\begin{proposition}\label{p:last}
The boredom series has a last member, and the last member is intentional.
\end{proposition}

\begin{proof}
Nothing that can be presented as dragging includes the time of one's existence as a proper part. A larger stretch of time, such as historical time or the time of the world, is not presented as dragging for me; it can be presented only as the time within which my existence falls, and it does not drag for anyone. So the third form is last. It is intentional by Lemma~4.3 and Theorem~4.4 of PHE-04-a.
\end{proof}

\begin{proposition}\label{p:noextension}
No conscious state lies beyond profound boredom in the boredom series.
\end{proposition}

\begin{proof}
A state beyond it would either present a larger object as dragging, which is excluded by Proposition~\ref{p:last}, or present no object at all while retaining the aspect. But an aspect is a way something is presented (Definition~2.1 of PHE-04-a), so a state with the aspect of dragging and no object would present nothing as dragging, and nothing is felt as dragging only if something is.
\end{proof}

Profound boredom is a limit case in two senses. It is the last member of a series in which the object widens, and it is the point at which the object ceases to be a being, since, as Lemma~4.2 of PHE-04-a shows, what profound boredom presents is not a being or an activity. What remains at the limit is time, and time is enough to secure intentionality.

\section{Two neighbors}

The result locates profound boredom relative to two neighbors. Anxiety, as PHE-04-a observes in Remark~4.6, presents not time but the nothing, and PHE-16 argues that what anxiety discloses is the withdrawal of beings as a whole, not a further entity. Boredom and anxiety are therefore limit cases along different dimensions: boredom along the dimension of time, anxiety along the dimension of beings. Both reach a point at which no being is the object, and in both a disclosure remains. The other neighbor is the functional account of boredom in EXI-10, on which boredom signals a mismatch between capacities and activities. The intentional analysis is presupposed by that account, and the limit case marks where the signal no longer concerns any particular activity; nothing here bears on whether, at that limit, the function is retained.

\section{Objections}

\begin{objection}
If profound boredom is intentional towards time, every awareness of time as passing would be boredom.
\end{objection}

\begin{reply}
The aspect matters as much as the object. Time can be presented as rushing, as full, as threatening; only time presented as dragging and empty is boredom. Definition~\ref{d:series} holds the aspect fixed for exactly this reason.
\end{reply}

\begin{objection}
Svendsen describes existential boredom as a lack of meaning in which nothing is presented at all \citep{Svendsen2005}. Calling its object \enquote{time} is an artefact of Heidegger's vocabulary.
\end{objection}

\begin{reply}
A lack of meaning is a lack in something: in one's life, which is lived as a stretch of time that is not filled. The description of the state as a lack therefore already gives it an object. What is missing in profound boredom is any being that matters, not anything presented.
\end{reply}

\section{Conclusion}

Profound boredom is the limit of a series in which the object of a constant aspect widens until it is the time of one's existence. The limit is intentional, and nothing lies beyond it in the series. The apparent exception to the intentionality thesis is its limit case, as PHE-04-a concludes.

\begin{thebibliography}{99}\small
\bibitem[Heidegger(1983)]{Heidegger1983} Heidegger, M. (1983). \emph{Die Grundbegriffe der Metaphysik. Welt -- Endlichkeit -- Einsamkeit}. Gesamtausgabe 29/30. Vittorio Klostermann.
\bibitem[Svendsen(2005)]{Svendsen2005} Svendsen, L. (2005). \emph{A Philosophy of Boredom}. Reaktion Books.
\end{thebibliography}

\end{document}
